\chapter{Untouchability: Forms and Perspectives}

\section{The Many Forms of Untouchability}

The syllabus says \textbf{'forms'} of untouchability because there is \textbf{no one form}. The common-sense idea is that in the caste system some people are untouchable --- but that is only common sense.

\begin{KeyConcept}
\begin{itemize}
\item The idea that 'some groups are untouchable' was objectified through the \textbf{Dharma Shastras} as interpreted by the Orientalists: they found a division between \textbf{Savarna} (those who have varna) and \textbf{Avarna} (those who have no varna --- hence untouchable).
\item Looked at through \textbf{spiritual rebirth}: the \textbf{Dvija} (twice-born --- Brahmin, Kshatriya, Vaishya) can be spiritually reborn, while the \textbf{Ekaj} (once-born --- Shudra and Avarna) cannot --- those \textbf{excluded from spirituality} were seen as untouchable (not allowed to read/interpret the Vedas or become saints).
\item Untouchability can be located \textbf{within Hinduism, across religions, within the Varna system (Shudra), across the Varna system (Avarna), across castes, within a caste, among upper-caste men and women, in the nation, and in gender} --- it is a \textbf{phenomenon}, not just a caste.
\end{itemize}
\end{KeyConcept}

\begin{KeyConcept}
\begin{itemize}
\item \textbf{Gandhi}: what the untouchables are to the upper-caste Hindu men, the same is \textbf{India to Britain} --- untouchability also cuts across to the category of \textbf{nation} (untouchable nation).
\item \textbf{Sabarimala} judgement: the social exclusion of women (based on menstruation) from the temple is \textbf{also a form of untouchability}.
\item Religious conversion did not end it: when low castes converted to Islam, Christianity or Buddhism, \textbf{their caste followed them} --- hence \textbf{Muslim Dalits, Dalit Christians, Dalit Buddhists}.
\end{itemize}
\end{KeyConcept}

\begin{ExampleBox}
\textbf{Satish Deshpande's reading of the astronaut joke} (10 candidates --- 5 OBC, 2 SC, 2 ST, 1 general astronaut): for the lower caste, \textbf{caste overrides qualification}; for the upper caste, \textbf{qualification overrides caste}. The victim was \textbf{Sameer Wankhede}, an IRS officer labelled a 'Muslim Dalit' --- the label itself shows we live in a Hindu upper-caste society.
\end{ExampleBox}

\section{Purity and Impurity as the Ideology (Dumont and A. M. Shah)}

\begin{KeyConcept}
\begin{itemize}
\item \textbf{Louis Dumont}: \textbf{purity/impurity is the ideology of Hindu civilization} --- not just of caste. It is the base of everything in Hindu life.
\item Everyday examples: removing slippers and washing feet before entering the sanctum; covering the head in a Gurdwara; the \textbf{sanctity of the kitchen} (the cooking mother lets no one in); restrictions on \textbf{food and water exchange} (food served by an untouchable becomes untouchable); endogamy in marriage; \textbf{childbirth} (the mother and child are secluded); \textbf{death} (sutak --- not visiting the bereaved home for about 12--13 days); not eating chicken during \textbf{Navratri}.
\item A \textbf{Jain} who observes sutak has, in Srinivas's terms, become \textbf{Sanskritized}.
\item \textbf{A. M. Shah} (Delhi School of Economics, following Dumont): \textbf{'Hindu civilization is sometimes called the civilization of purity and pollution,'} and the Hindu psyche is believed to be \textbf{pathologically obsessed} with them.
\end{itemize}
\end{KeyConcept}

\begin{ExampleBox}
\begin{itemize}
\item A. M. Shah's illustration (a Gujarati novel, \textbf{Neelkanth}, c. 1900): a Brahmin named \textbf{Bhadram} travelling from Ahmedabad to Bombay by train considered his co-passengers polluting and took a \textbf{purificatory bath at every station}.
\item The word \textbf{asprishya} (the Sanskrit equivalent of 'untouchable') was used \textbf{even for the highest purity}: in the \textbf{Pushti Marg} (a Vaishnava sect founded by \textbf{Vallabhacharya}), a worshipper in a state of intense purity --- e.g. \textbf{Meera Bai} worshipping Krishna --- was in an 'asprishya state', which others could not touch or disturb. So untouchability applies to the \textbf{most impure and the most pure alike}.
\end{itemize}
\end{ExampleBox}

\section{Temporary vs Permanent Pollution}

\begin{KeyConcept}
\begin{itemize}
\item Dumont's \textbf{two types of impurity}:
\item \textbf{Permanent pollution} --- the \textbf{low castes}, involved in \textbf{menial occupations} (manual scavenging, skinning animals), seen as permanently impure.
\item \textbf{Temporary pollution} --- the \textbf{upper castes}: death (sutak), childbirth, menstruation --- impurity that can be \textbf{purified through rituals}.
\item Even \textbf{upper-caste women} are the victims of temporary impurity --- the biggest being the \textbf{menstrual cycle} (during which marriage is not even discussed) and \textbf{childbirth} (mother and child are secluded).
\end{itemize}
\end{KeyConcept}

\section{Untouchability Within Brahmins and Among Non-Brahmins}

\begin{KeyConcept}
\begin{itemize}
\item \textbf{Brahmins are not a homogeneous group} --- so there are \textbf{untouchable Brahmins}. Among Brahmins: the \textbf{temple priest} is the purest; the \textbf{ordinary} Brahmin is in the middle; the Brahmin specialising in \textbf{mortuary (death) rituals} is the lowest --- a death-ritual priest is not called for marriages.
\item \textbf{Beteille}: Brahmins are divided into segments, each a \textbf{status group} with its own lifestyle --- so the Smartha do not marry the Sri Vaishnava (they treat each other as untouchable), and purity is reproduced generation after generation within each segment.
\item \textbf{Non-Brahmins} of highly Sanskritic sects (e.g. the \textbf{Bhagats} and \textbf{Bhuvas} --- low-caste/tribal sects) observe purity rules so meticulously that they consider ordinary Brahmins \textbf{less pure}, even polluting.
\item Among the \textbf{Devadasis} there is also a hierarchy: the \textbf{Deo Dasi} (temple) is purest, the \textbf{domestic Dasi} is most impure. (The scholar \textbf{Subhadra Mitra Channa} speaks of Hindu women's two statuses --- \textbf{Devi and Dasi}.)
\end{itemize}
\end{KeyConcept}

\section{The Career of the Concept: From Avarna to Scheduled Caste}

\begin{KeyConcept}
\begin{itemize}
\item The word \textbf{untouchable} is a modern (English) category; the traditional Sanskrit category was \textbf{Avarna} --- the precursor of the modern 'untouchables'.
\item The concept has had its own \textbf{career} (Simon Charsley's essay): \textbf{Avarna $\rightarrow$ Asprishya Shudra $\rightarrow$ Depressed Classes $\rightarrow$ Excluded Classes $\rightarrow$ Harijan $\rightarrow$ Scheduled Caste $\rightarrow$ Dalit}.
\item \textbf{Herbert Risley}, the 1901 census commissioner, was the first to classify the Shudra varna into five categories, including \textbf{Asprishya Shudra} ('untouchable Shudra') --- but no coherent, consistent list ever emerged.
\item \textbf{Jyotiba Phule} first spoke of the \textbf{Dalit} category; \textbf{Ambedkar} used 'Depressed Classes'; \textbf{Gandhi} gave 'Harijan'; the Government of India gave \textbf{Scheduled Caste}.
\item In preparing these lists, officials always agreed on the \textbf{scavengers} (the lowest) but disagreed on the others --- because there was \textbf{no clear criterion}. As a concept, 'untouchable' \textbf{suppressed the diversity} of these groups.
\end{itemize}
\end{KeyConcept}

\begin{KeyConcept}
\textbf{Michel Foucault}: knowledge constitutes power --- the people who manufactured this knowledge (of who is 'untouchable') acquired power over our minds; the 'fixed, inviolable line' between touchable and untouchable is a construction, not objective reality.
\end{KeyConcept}

\section{Untouchables Were Never Homogeneous}

\begin{KeyConcept}
\begin{itemize}
\item The ancient texts do \textbf{not} mention the numerous jatis of untouchables --- only \textbf{field studies} reveal that untouchables are themselves divided into endogamous castes arranged in hierarchy (some 'pure', some 'impure'), just like the rest of society.
\item The \textbf{Smiths of Rampura} claimed \textbf{Vishwakarma Brahmin} status.
\item \textbf{Garo} (Gujarat) --- an untouchable \textbf{priest} caste that claims Brahmin status, wears the sacred thread, and adopts Brahminical surnames (Joshi, Vyas) --- pure Sanskritization. They form part of a \textbf{hierarchy of priestly groups}: the high Brahmin sub-castes served as priests for high non-Brahmin castes, while a few low Brahmin sub-castes served lower-middle castes --- status matters more than power (a Brahmin priest would not serve even a wealthy Dalit).
\end{itemize}
\end{KeyConcept}

\begin{ExampleBox}
\textbf{Untouchable saints}: the book \textbf{'Untouchable Saints'} (by \textbf{Eleanor Zelliot} and \textbf{Rohini Mokashi-Punekar}, 2005) shows that since about the 10th century there were untouchable saint-poets --- \textbf{Tiruppan} (given the high \textbf{Alvar} status) and \textbf{Nandanar} (given \textbf{Nayanar} status) in Tamil Nadu, \textbf{Chokha Mela} in Maharashtra, and \textbf{Ravidas} in North India --- whose poetry was accepted by high castes into temple and domestic worship. The \textbf{Bhakti movement} was a challenge to mainstream Brahminical philosophy, led by the low castes.
\end{ExampleBox}

\section{Impact of Modernity: The Transfer of Purity}

\begin{KeyConcept}
\begin{itemize}
\item Ideas like \textbf{humanitarianism, meritocracy, liberty, fraternity, fundamental rights, urbanization, industrialization, secularization and democracy} diluted purity/impurity --- but \textbf{only among the westernized upper-caste Hindu men}.
\item Meanwhile the \textbf{lower castes were getting Sanskritized} --- absorbing the very purity/pollution ideas that the upper castes were giving up. So \textbf{purity/impurity did not die; it was transferred from the upper castes to the lower castes} (the status quo was maintained).
\item \textbf{A. M. Shah's summary}: in \textbf{pre-modern India}, concern for purity/pollution \textbf{decreased as one went down the hierarchy}; in \textbf{modern India} there is a \textbf{two-way change} --- decreasing among the upper castes (due to westernization) and increasing among the lower castes (due to Sanskritization).
\end{itemize}
\end{KeyConcept}

\begin{QuickRevision}
\begin{itemize}
\item Forms of untouchability: within/across Hinduism, within/across Varna, across castes, within caste, upper-caste men \& women, nation, gender --- a phenomenon, not just caste.
\item Savarna vs Avarna; Dvija (twice-born) vs Ekaj (once-born, excluded from spirituality).
\item Dumont: purity/impurity = ideology of Hindu civilization; A. M. Shah: 'civilization of purity and pollution'.
\item Temporary (upper caste: death/birth/menstruation) vs permanent (low caste: menial occupations) pollution.
\item Within Brahmins: temple priest > ordinary > mortuary-ritual Brahmin (Beteille's segments); non-Brahmin Sanskritic sects (Bhagats, Bhuvas); Devadasi (Deo Dasi vs domestic Dasi).
\item Career of the concept: Avarna $\rightarrow$ Asprishya Shudra (Risley 1901) $\rightarrow$ Depressed Classes (Ambedkar) $\rightarrow$ Harijan (Gandhi) $\rightarrow$ Scheduled Caste; no clear criterion (only scavengers agreed); suppressed diversity; Foucault (knowledge = power).
\item Untouchables never homogeneous (field studies): Smiths of Rampura (Vishwakarma Brahmin), Garo of Gujarat (Brahmin claim); untouchable saints (Tiruppan/Alvar, Nandanar/Nayanar, Chokha Mela, Ravidas; Zelliot \& Mokashi-Punekar).
\item Impact of modernity: purity/impurity transferred from (westernized) upper castes to (Sanskritized) lower castes --- two-way change.
\end{itemize}
\end{QuickRevision}

